% !TeX root = ../../main.tex
\subsection{Signed loop moments}
\label{subsec:s04:01}

An edge has one integration parameter, shared by its two flags in
\eqref{eq:s03:integrated-graph-formula}. For nonnegative integers
$a,b$ define its paired self-loop weight by
\begin{equation}\label{eq:s04:paired-weight}
 w_{ab}=\frac12\int_0^\alpha
       \frac{(-t/2)^{a+b}}{(2a+1)!(2b+1)!}\,\dd t
 =\frac{(-1)^{a+b}\pi^{2a+2b+2}}
       {(a+b+1)(2a+1)!(2b+1)!}.
\end{equation}
The factor $1/2$ accounts for exchange of the two branches of the loop. The denominator
$a+b+1$ records their common parameter; their weights do not factor
into independent endpoint weights.

For any nonnegative integer degree tuple $\mathbf d$, write
$D=|\mathbf d|$ and let
$m_a=\#\{i:d_i=a\}$ for $a=0,1,2$. Given $p\geq0$ ordered pairs
$(a_i,b_i)$, use the concatenated tuple
$\mathbf d=(a_1,b_1,\ldots,a_p,b_p)$, of length $2p$.
For a polynomial $\Phi$ in $p,D,m_0,m_1,m_2$ set
\begin{equation}\label{eq:s04:loop-functional}
 \mathcal L[\Phi]=e^{-S(1)}\sum_{p\geq0}\frac1{p!}
 \sum_{a_1,b_1,\ldots,a_p,b_p\geq0}
   \left(\prod_{i=1}^p w_{a_i b_i}\right)
       \Phi(p,D,m_0,m_1,m_2).
\end{equation}
The empty tuple has all five statistics zero. Together these symmetry
factors give $1/(2^pp!)$, exactly the automorphism factor for the
restored root loops. We also use the same series for the exponential
test functions in the next lemma.

The joint generator retains the low degrees needed at second order.
Let $v,x,z_0,z_1,z_2$ be independent complex moment parameters.

\begin{lmm}[Joint loop moments]\label{lmm:s04:loop-moments}
The functional~\eqref{eq:s04:loop-functional} is absolutely convergent
on polynomials, satisfies $\mathcal L[1]=1$, and has the entire
joint generating function
\begin{align}
 &\mathcal L\left[e^{vp+xD+\sum_{a=0}^2z_a m_a}\right]\notag\\
 &\quad=\exp\left\{\frac{e^v}{2}\int_0^\alpha
 \left[f(e^x t)+\sum_{a=0}^2(e^{z_a}-1)
               \frac{(-e^x t/2)^a}{(2a+1)!}\right]^2\dd t-S(1)\right\}.
 \label{eq:s04:joint-generator}
\end{align}
Every fixed moment is obtained by differentiating at zero. The
functional is signed, rather than a probability expectation.
\end{lmm}
\begin{prf}
For one edge the square selects its two endpoint degrees. Its
integral gives $w_{ab}$, the factors $e^x$ count their total degree,
and $e^{z_a}$ counts endpoints of degree $a$. Summing over $p$ with
$1/p!$ exponentiates this one-edge sum.

On compact sets of the moment parameters, the absolute endpoint
series are bounded by factorial series on the finite interval
$[0,\alpha]$. Their exponential sum over $p$ converges normally,
also after any fixed number of derivatives. This proves the
convergence and differentiation assertions. At zero parameters the
one-edge term is $\tfrac12\int_0^\alpha f(t)^2\dd t=S(1)$, proving
$\mathcal L[1]=1$.
\end{prf}

\begin{exm}[Two moments]
For a local complex variable $z$ near $1$, the degree generator of
one loop is
\[
 \sum_{a,b\geq0}w_{ab}z^{a+b}
 =\frac12\int_0^\alpha f(zt)^2\dd t=\frac{S(z)}z.
\]
Taking its Euler derivative gives the corresponding logarithmic
derivative of~\eqref{eq:s04:joint-generator}. Since
$S'(1)=\alpha f(\alpha)^2/2=0$, it follows that
\begin{equation}\label{eq:s04:representative-moments}
 \mathcal L[D]=-S(1),\qquad \mathcal L[m_0]=8.
\end{equation}
For the second identity, differentiation in $z_0$ gives
$\int_0^\alpha f(t)\dd t$. With $t=2y^2$ this is
$4\int_0^\pi\sin y\dd y=8$. Both endpoints contribute to this
derivative of the square. The remaining finite differentiation rules
and moment tables are given in Subsection~\ref{subsec:a:02}.
\end{exm}
